\documentclass[10pt,a4paper,twocolumn]{article}

\usepackage[brazil]{babel}
\usepackage[utf8]{inputenc}
\usepackage[T1]{fontenc}
\usepackage[left=1.6cm,right=1.6cm,top=1.7cm,bottom=1.7cm]{geometry}
\usepackage{graphicx}
\usepackage{amsmath,amssymb}
\usepackage{booktabs,multirow,array}
\usepackage[hidelinks]{hyperref}
\usepackage{fancyhdr}
\usepackage{caption}

\pagestyle{fancy}
\fancyhf{}
\lhead{Modelo de Artigo em Duas Colunas}
\rhead{\thepage}
\renewcommand{\headrulewidth}{0.4pt}

\title{\textbf{Título do Artigo em Duas Colunas}}
\author{
Nome do Autor \\
\small Universidade / Instituição \\
\small \texttt{email@exemplo.com}
}
\date{}

\begin{document}
\maketitle

\begin{abstract}
Este modelo foi criado para servir como base de artigos acadêmicos em duas colunas, com estrutura simples, clara e adaptável. O objetivo é demonstrar organização textual, uso de seções, figuras, tabelas e formatação adequada para trabalhos técnicos.
\end{abstract}

\noindent\textbf{Palavras-chave:} artigo acadêmico; duas colunas; \LaTeX; diagramação.

\section{Introdução}
Este texto apresenta um modelo de artigo em duas colunas. Ele pode ser usado como base para trabalhos acadêmicos, relatórios de pesquisa e manuscritos técnicos.

\section{Estrutura do documento}
A estrutura foi pensada para facilitar a leitura e a organização do conteúdo. Cada seção pode ser expandida conforme a necessidade do autor.

\subsection{Exemplo de referência cruzada}
Como demonstrado na Seção~\ref{sec:metodologia}, a divisão do trabalho em etapas melhora a clareza do manuscrito.

\section{Metodologia}
\label{sec:metodologia}
A metodologia pode incluir pesquisa bibliográfica, desenvolvimento experimental, análise de dados ou estudo de caso. O importante é manter consistência entre objetivos, métodos e resultados.

\section{Exemplo de figura}
\begin{figure}[t]
    \centering
    \fbox{\rule{0pt}{3.5cm}\rule{0.9\linewidth}{0pt}}
    \caption{Figura de exemplo com placeholder}
    \label{fig:exemplo}
\end{figure}

\section{Exemplo de tabela}
\begin{table}[t]
\centering
\caption{Tabela resumida de exemplo}
\label{tab:exemplo}
\begin{tabular}{lcc}
\toprule
Item & Valor A & Valor B \\
\midrule
Linha 1 & 10 & 12 \\
Linha 2 & 15 & 18 \\
Linha 3 & 20 & 25 \\
\bottomrule
\end{tabular}
\end{table}

\section{Resultados}
Os resultados devem ser apresentados de forma objetiva, com apoio de tabelas, figuras e comparações quando necessário.

\section{Conclusão}
Este modelo é adequado para documentos simples e claros em duas colunas, oferecendo uma base profissional para publicações acadêmicas.

\end{document}